\SourcePage{224}{229}
\section{Radiation of diatomic molecules. Electronic spectra}

The specificity of molecular spectra is connected first of all with the separation of the energy of the molecule into electronic, vibrational and rotational parts, of which each succeeding one is small compared with the preceding. The structure of levels of diatomic molecules was studied in detail in vol.~III, Ch.~XI. Here we shall concern ourselves with clarifying the picture of the spectrum that arises and with computing the intensities of lines in it\,\footnote{The subsequent exposition is based on the material in vol.~III, \S\S\,78, 82--88. To avoid cluttering the text, we shall refrain from constant references to these paragraphs.}.

Let us start with the general case, when in a transition the electronic state of the molecule changes (and along with it, generally speaking, the vibrational and rotational states as well). The frequencies of these transitions lie in the visible and ultraviolet regions of the spectrum. The totality of them is spoken of as the \textit{electronic spectrum} of the molecule. We shall always have in mind electric-dipole transitions; transitions of other types are generally of little significance in molecular spectroscopy.

As for dipole transitions in any system, the selection rule for the total angular momentum $J$ of the molecule holds:
\begin{equation}
|J' - J| \leqslant 1 \leqslant J + J'.
\tag{53.1}
\end{equation}

The strict selection rule for the parity of the system corresponds in this case to a selection rule for the \textit{sign} of the level; recall that in the molecular spectroscopy terminology adopted, states with wave functions not changing or changing sign on inversion (change of sign of the coordinates of electrons and nuclei) are called \textit{positive} or \textit{negative}. Thus, we have the strict rule:
\begin{equation}
+ \to -,\qquad - \to +.
\tag{53.2}
\end{equation}

If the molecule consists of identical atoms (with nuclei of one and the same isotope), there arises the classification of levels according to the behaviour of the wave functions under permutation of the nuclear coordinates: \textit{symmetric} ($s$) levels with wave functions not changing sign under this transformation, and \textit{antisymmetric} ($a$) levels with functions changing sign. Since the operator of the electronic dipole moment is not affected by this transformation at all, its matrix elements are different from zero only for transitions without change of this symmetry\,\footnote{This rule applies, obviously, also to transitions of any multipolarity.}:
\begin{equation}
s \to s,\qquad a \to a.
\tag{53.3}
\end{equation}

\SourcePage{225}{230}
This rule, however, is not absolutely strict. The point is that the existence in levels of a given symmetry property is connected with the existence in the molecule of one or another definite value of the total nuclear spin $I$. By virtue of the extreme weakness of the interaction of the nuclear spins with the electrons, $I$ is conserved with great accuracy, but still not strictly. When this interaction is taken into account, $I$ will not have a definite value, the symmetry property ($s$ or $a$) is not conserved, and the selection rule~(53.3) is lost.

The electronic terms of molecules of identical atoms are characterised also by their \textit{parity} ($g$ or $u$)---the behaviour of the wave functions (as measured from the centre of the molecule) under change of sign of the electron coordinates at unchanged nuclear coordinates. There exists a close connection between this property of the electronic term, on the one hand, and the nuclear symmetry and the sign of the rotational levels belonging to this term, on the other. The levels belonging to an even ($g$) electronic term may have characteristics $s+$ or $a-$, and those to an odd one $s-$ or $a+$. From the rules~(53.2) and~(53.3) there follows therefore also the rule
\begin{equation}
\mathrm{g} \to \mathrm{u},\qquad \mathrm{u} \to \mathrm{g}.
\tag{53.4}
\end{equation}

As an approximate rule~(53.4) remains valid also for molecules from different isotopes of one and the same element. Since the nuclear charges are identical, on considering the electronic term with the nuclei at rest we are dealing in such a case with a system of electrons in an electric field possessing a centre of symmetry (at the point dividing in half the distance between the nuclei). The symmetry of the electronic wave function with respect to inversion at this point determines the parity of the term, and since the vector of the electric dipole moment changes sign under this transformation, the parity must change in a transition; hence we arrive at the rule~(53.4). The approximateness of the rule, based on this derivation alone, is connected with the necessity of treating the nuclei as stationary. It therefore breaks down when the interaction between the electronic state and the rotation of the molecule is taken into account.

Further selection rules are connected with one or another concrete assumptions about the comparative magnitudes of the various interactions in the molecule (i.e.\ about its coupling type). These rules can therefore be only approximate.

The majority of electronic terms of diatomic molecules belong to coupling types $a$ or $b$. Both these types are characterised by the fact that the coupling of the orbital angular momentum with the axis of the molecule (the electric interaction of both atoms in the molecule) is large compared with all other interactions. In this connection there exist

\SourcePage{226}{231}
the quantum numbers $\Lambda$ and $S$ (the projection of the electronic orbital angular momentum onto the axis of the molecule and the total spin of the electrons). The operator of the orbital quantity---the electronic orbital angular momentum---commutes with the spin operator, so that
\begin{equation}
S' - S = 0\qquad (\text{case } a,\,b).
\tag{53.5}
\end{equation}
The number $\Lambda$ is subject to the selection rule
\begin{equation}
\Lambda' - \Lambda = 0,\,\pm 1\qquad (\text{case } a,\,b),
\tag{53.6}
\end{equation}
and for transitions between states with $\Lambda = 0$ ($\Sigma$-terms) the additional rule applies
\begin{equation}
\Sigma^+ \to \Sigma^+,\qquad \Sigma^- \to \Sigma^-\qquad (\text{case } a,\,b)
\tag{53.7}
\end{equation}
(recall that the states $\Sigma^+$ and $\Sigma^-$ are distinguished by their behaviour under reflection in a plane passing through the axis of the molecule). The rules~(53.6), (53.7) are obtained in the system of coordinates rigidly connected with the nuclei (see~III, \S\,87); the rule~(53.6) is analogous to the selection rule for the magnetic quantum number in the atomic case.

The coupling types $a$ and $b$ differ from each other by the ratio between the energy of the ``spin--axis'' interaction and the energy of rotation (the spacings of the rotational levels). In case $a$ the first of these is larger, while in case $b$ the other way round.

\textbf{Case $a$.} In this case there exists the quantum number $\Sigma$---the projection of the total spin onto the axis of the molecule (and with it the number $\Omega = \Sigma + \Lambda$---the projection of the total angular momentum). If both (initial and final) states belong to case $a$, the rule
\begin{equation}
\Sigma' - \Sigma = 0\qquad (\text{case } a)
\tag{53.8}
\end{equation}
is valid (following from the already mentioned commutativity of the dipole moment operator with the spin). From~(53.6) and~(53.8) follows
\begin{equation}
\Omega' - \Omega = 0,\,\pm 1.
\tag{53.9}
\end{equation}

If $\Omega = \Omega' = 0$, then additionally to the general rule~(53.1) transitions with $J' = J$ are forbidden\,\footnote{This rule is analogous to the prohibition of transitions with $J = J'$ at $M = M' = 0$ for atoms (see~\S\,5), where, however, it could be of interest only in the presence of an external field. In the present case the rule follows directly from the formula~(53.12) given below; the $3j$-symbol $\begin{pmatrix} J' & 1 & J \\ 0 & 0 & 0\end{pmatrix}$ vanishes at $J' = J$, when the sum $J' + J + 1$ is odd.}:
\begin{equation}
J' - J = \pm 1\quad\text{at}\quad \Omega = \Omega' = 0\qquad (\text{case } a).
\tag{53.10}
\end{equation}

\SourcePage{227}{232}
Let us consider transitions between the rotational components of any two definite vibrational levels belonging to two different electronic terms (of type $a$). Taking into account the fine structure of the electronic term, each of these levels splits into several components, the number of which $(2S + 1)$ must by virtue of rule~(53.5) be the same in both. According to rule~(53.8) each component of one level combines only with a single component of the other with the same value of $\Sigma$.

Let us take further a pair of levels with identical $\Sigma$; their values of $\Omega$ and $\Omega'$ can differ (together with $\Lambda$ and $\Lambda'$) by 0 or $\pm 1$. When the rotation is taken into account each of them splits into a sequence of levels, distinguished by consecutive values of $J$ and $J'$ running from values $J \geqslant |\Omega|$, $J' \geqslant |\Omega'|$. The dependence of the transition probabilities on these numbers can be established in the general form (\textit{H.\,H\"onl, F.\,London}, 1925).

The matrix element of the transition $n\Lambda\Omega J M_J \to n'\Lambda'\Omega'J'M'_J$ ($n$ denoting the remaining quantum numbers besides $\Omega$ and $\Lambda$, characteristics of the electronic term) equals
\begin{align}
|\langle n'\Lambda'\Omega'J'M'_J|d_q|n\Lambda\Omega JM_J\rangle| &={}\notag\\[4pt]
= \sqrt{(2J+1)(2J'+1)}&\begin{pmatrix} J' & 1 & J \\ -\Omega' & q' & \Omega\end{pmatrix}\begin{pmatrix} J' & 1 & J \\ -M'_J & q & M_J\end{pmatrix}\times{}\notag\\[4pt]
&\times|\langle n'\Lambda'|\bar{d}_{q'}|n\Lambda\rangle|,
\tag{53.11}
\end{align}
where $d_q$ and $\bar{d}_{q'}$ are the spherical components of the vector of the dipole moment respectively in the fixed system of coordinates $xyz$ and in the ``moving'' system $\xi\eta\zeta$ with the axis $\zeta$ along the axis of the molecule (this formula is obtained with the help of formula~(110.6) (see~III)). The matrix elements $\langle n'\Lambda'|\bar{d}_{q'}|n\Lambda\rangle$ do not depend on the rotational quantum numbers $J$, $J'$, but depend only on the characteristics of the electronic terms\,\footnote{One can verify this in the same way as was done at the beginning of~III, \S\,29 for a scalar quantity $f$. In the present case the operator of the vectorial quantity $\mathbf{d}$ commutes with the operator of the conserved (in the zeroth approximation) vector $\mathbf{S}$ on the axis $\boldsymbol{\zeta}$, and $\Sigma$ is the projection of $\mathbf{S}$ onto the axis of the rotating system of coordinates, in which one has to consider the condition of commutativity of $\mathbf{d}$ and $\mathbf{S}$.}; therefore in the designation of the matrix element the indices $\Omega$ and $\Omega'$ are omitted.

The probability of the transition $n\Lambda\Omega J \to n'\Lambda'\Omega'J'$ is proportional to the square of the matrix element~(53.11), summed over $M'_J$. By virtue of formula~(106.12) (see~III) we have
\[
\sum_{M'_J}\begin{pmatrix} J' & 1 & J \\ -M'_J & q & M_J\end{pmatrix}^{\!2} = \frac{1}{2J+1},
\]

\SourcePage{228}{233}
so that
\begin{equation}
w(n\Lambda\Omega J \to n'\Lambda'\Omega'J') = (2J'+1)\begin{pmatrix} J' & 1 & J \\ -\Omega' & \Omega'-\Omega & \Omega\end{pmatrix}^{\!2}B(n',n;\Lambda',\Lambda),
\tag{53.12}
\end{equation}
where the coefficients $B$ do not depend on $J$, $J'$ (we neglect, of course, the small difference in frequencies of transitions between different $J$, $J'$)\,\footnote{If we sum~(53.12) over $J'$, then (by the orthogonality property of the $3j$-symbols (see~III, (106.13))) we get simply $B(n',n;\,\Lambda',\Lambda)$. In other words, the total probability of the transition from the rotational level $J$ of the state $\Omega$ to all levels $J'$ of the state $\Omega'$ does not depend on~$J$.}.

\textbf{Case $b$.} In this case there exists, alongside the total angular momentum $J$, also the quantum number $K$---the angular momentum of the molecule without the spin. The selection rules for this number coincide with the general selection rules for any orbital vectorial quantity (which the electric dipole moment is):
\begin{equation}
|K' - K| \leqslant 1 \leqslant K + K'\qquad (\text{case } b)
\tag{53.13}
\end{equation}
with the additional prohibition of the transition with $K = K'$ at $\Lambda = \Lambda' = 0$ (analogously to~(53.10)):
\begin{equation}
K' - K = \pm 1\quad\text{at}\quad \Lambda = \Lambda' = 0.
\tag{53.14}
\end{equation}

Let us consider the transitions between the rotational components of definite vibrational levels of two electronic states belonging to type $b$. The probabilities of transitions between them are determined by the same formulae~(53.12), in which one has to write $K$, $\Lambda$ instead of $J$, $\Omega$. When the fine structure is taken into account (for $S \neq 0$) each rotational level $K$ splits into $2S + 1$ components with $J = |K - S|$, \ldots, $K + S$, as a result of which instead of a single line $J \to J'$ a multiplet arises. Since in the present case we are dealing with the composition of the free (not connected with the molecular axis) angular momenta $\mathbf{K}$ and $\mathbf{S}$, the formulae for the relative probabilities of the various lines of the multiplet coincide with the analogous formulae~(49.15) for the fine-structure components of atomic spectral lines, where the same role is played by the angular momenta $\mathbf{L}$ and $\mathbf{S}$ in the case of $LS$-coupling\,\footnote{Each of the considered rotational levels $J$ when the $\Lambda$-doubling is taken into account splits further into two levels, of which one is positive and the other negative. Therefore instead of a single transition $J \to J'$ one must take into account, taking into account the selection rule~(53.2), two transitions: from the positive (negative) component of level $J$ to the negative (positive) component of level $J'$. The probabilities of these transitions are equal.}.

\SourcePage{229}{234}
Thus, we have examined the selection rules determining the possible lines of the spectrum in all the main cases that can arise in diatomic molecules.

The set of lines from transitions between rotational components of two given electronic-vibrational levels constitutes, as one says in spectroscopy, a \textit{band}; by virtue of the smallness of the rotational intervals the lines in the band are very closely spaced. The frequencies of these lines are given by the expressions
\begin{equation}
\hbar\omega_{J,J'} = \mathrm{const} + BJ(J+1) - B'J'(J'+1),
\tag{53.15}
\end{equation}
where $B$, $B'$ are the rotational constants in both electronic states (to avoid superfluous complications we assume the electronic terms to be singlet). For $J' = J$, $J \pm 1$ formula~(53.15) is represented graphically (Fig.~2) by three branches (parabolas), whose points for integer values of $J$ give the frequencies. (The arrangement of branches in Fig.~2 corresponds to the case $B' < B$; for $B' > B$ they are opened in the direction of small $\omega$, and the upper branch is the curve for $J' = J - 1$\,\footnote{Series of lines corresponding to transitions with $J' = J + 1$, $J$, $J - 1$ are called respectively $P$-, $Q$- and $R$-branches.}.) The presence of a turning branch leads, as is clear from the figure, to a bunching of lines towards a definite limiting position (\textit{band edges}).

Speaking of intensities of lines, one should also mention a peculiar phenomenon of \textit{intensity alternation} in certain bands of the electronic spectrum of molecules of atoms of one and the same isotope (\textit{W.\,Heisenberg, F.\,Hund}, 1927). The symmetry requirements connected with the nuclear spins lead to the fact that in the electronic $\Sigma$-terms the rotational components with even and odd values of $K$ possess opposite symmetry with respect to the nuclei and correspondingly different nuclear statistical weights. According to rule~(53.14), when the transitions between two different $\Sigma$-terms are considered, by virtue of rule~(53.4) one of the $\Sigma$-terms must be even and the other odd. As a result, for a given value of $J' - J$ the transitions alternate between pairs of symmetric and pairs of antisymmetric levels (as is illustrated by the diagram in Fig.~3 for the case of the states $\Sigma^+_u$ and $\Sigma^+_g$). On the other hand, the observed intensity of a line is proportional to the number of molecules found in a given initial state, and thereby to its statistical weight. Therefore the intensity of consecutive lines ($J = 0$, 1, 2, \ldots) will alternate between larger and smaller, proportional alternately to $g_a$ and

\SourcePage{230}{235}
$g_s$ (besides the monotonic variation predictable from formulae~(53.12))\,\footnote{It is assumed that all states with different values of the total nuclear spin are populated uniformly.}.

\begin{figure}[h]
\centering
\setlength{\unitlength}{1mm}
\begin{picture}(50,32)
\put(5,0){\vector(0,1){30}}
\put(5,0){\vector(1,0){45}}
\put(50,0){\makebox(0,0)[l]{\footnotesize $\omega$}}
\put(0,30){\makebox(0,0)[r]{\footnotesize $J$}}
\put(10,2){\qbezier(0,0)(10,14)(25,27)}
\put(20,2){\qbezier(0,0)(8,11)(18,26)}
\put(28,2){\qbezier(0,0)(6,8)(14,25)}
\put(38,27){\makebox(0,0)[l]{\footnotesize $J' = J+1$}}
\put(40,24){\makebox(0,0)[l]{\footnotesize $J' = J$}}
\put(44,21){\makebox(0,0)[l]{\footnotesize $J' = J-1$}}
\end{picture}
\hspace{1cm}
\begin{picture}(50,32)
\put(5,0){\vector(0,1){30}}
\put(5,0){\vector(1,0){45}}
\put(50,0){\makebox(0,0)[l]{\footnotesize $r$}}
\put(15,2){\qbezier(0,0)(0,13)(0,25)}
\put(35,2){\qbezier(0,0)(0,13)(0,25)}
\put(10,14){\makebox(0,0)[r]{\footnotesize $U(r)$}}
\put(40,14){\makebox(0,0)[l]{\footnotesize $U'(r)$}}
\put(10,22){\line(1,0){20}}
\put(10,10){\line(1,0){20}}
\put(32,22){\makebox(0,0)[l]{\footnotesize $E$}}
\put(32,10){\makebox(0,0)[l]{\footnotesize $E'$}}
\end{picture}

\vspace{2mm}
\makebox[50mm]{\footnotesize Fig.~2}\hspace{1cm}\makebox[50mm]{\footnotesize Fig.~4}
\end{figure}

\begin{figure}[h]
\centering
\setlength{\unitlength}{1mm}
\begin{picture}(80,35)
\put(0,30){\makebox(0,0)[r]{\footnotesize $\Sigma^+_u$}}
\put(0,5){\makebox(0,0)[r]{\footnotesize $\Sigma^+_g$}}
\multiput(5,30)(12,0){6}{\line(1,0){8}}
\multiput(5,5)(12,0){6}{\line(1,0){8}}
\put(5,32){\makebox(0,0){\footnotesize $J=0$}}
\put(17,32){\makebox(0,0){\footnotesize 1}}
\put(29,32){\makebox(0,0){\footnotesize 2}}
\put(41,32){\makebox(0,0){\footnotesize 3}}
\put(53,32){\makebox(0,0){\footnotesize 4}}
\put(65,32){\makebox(0,0){\footnotesize 5}}
\put(5,27){\makebox(0,0){\footnotesize $a$}}
\put(17,27){\makebox(0,0){\footnotesize $s$}}
\put(29,27){\makebox(0,0){\footnotesize $a$}}
\put(41,27){\makebox(0,0){\footnotesize $s$}}
\put(53,27){\makebox(0,0){\footnotesize $a$}}
\put(65,27){\makebox(0,0){\footnotesize $s$}}
\put(5,2){\makebox(0,0){\footnotesize $J'=0$}}
\put(17,2){\makebox(0,0){\footnotesize 1}}
\put(29,2){\makebox(0,0){\footnotesize 2}}
\put(41,2){\makebox(0,0){\footnotesize 3}}
\put(53,2){\makebox(0,0){\footnotesize 4}}
\put(65,2){\makebox(0,0){\footnotesize 5}}
\put(5,8){\makebox(0,0){\footnotesize $s$}}
\put(17,8){\makebox(0,0){\footnotesize $a$}}
\put(29,8){\makebox(0,0){\footnotesize $s$}}
\put(41,8){\makebox(0,0){\footnotesize $a$}}
\put(53,8){\makebox(0,0){\footnotesize $s$}}
\put(65,8){\makebox(0,0){\footnotesize $a$}}
\put(9,29){\line(1,-2){10.5}}
\put(21,29){\line(1,-2){10.5}}
\put(33,29){\line(1,-2){10.5}}
\put(45,29){\line(1,-2){10.5}}
\put(57,29){\line(1,-2){10.5}}
\end{picture}

\vspace{2mm}
\makebox[80mm]{\footnotesize Fig.~3}
\end{figure}

For changes in the vibrational quantum number in transitions between two different electronic terms there are no strict selection rules. There exists, however, a rule (the \textit{Franck--Condon principle}), allowing one to predict the most probable change of the vibrational state. It is based on the quasiclassical character of the motion of the nuclei, connected with their large mass (cf.\ what was said about predissociation in~III, \S\,90)\,\footnote{Strictly speaking, it is also necessary that the vibrational quantum number be sufficiently large.}.

In the integral determining the matrix element of the transition between vibrational states with energies $E$ and $E'$ of electronic terms $U(r)$ and $U'(r)$, the main role is played by the neighbourhood of the point $r = r_0$, at which
\begin{equation}
U(r_0) - U'(r_0) = E - E'
\tag{53.16}
\end{equation}
(i.e.\ the momenta of the relative motion of the nuclei in both states are the same: $p = p'$). For a given value of $E$ the probability of the transition (as a function of the final energy $E'$) is the greater the smaller each of the differences $E - U$ and $E' - U'$. It is maximal when
\begin{equation}
E - U(r_0) = E' - U'(r_0) = 0,
\tag{53.17}
\end{equation}
i.e.\ when the ``transition point'' $r_0$ (the root of equation~(53.16)) coincides with the classical turning point of the nuclei (Fig.~4 illustrates this connection between $E$ and the most probable $E'$ graphically). For clarity one can say that the most probable transition occurs near the point at which the nuclei come to rest and in the neighbourhood of which they spend, consequently, the most time.
